\section*{Capítulo \ref{ch:b3:point-groups} --- Simetria e grupos pontuais}

\begin{solution}{exo:b3:point-groups:1}
Eteno: três $C_2$ perpendiculares, três planos, $i$: $D_{2h}$. \ce{XeF4} (quadrado
plano): $C_4$, quatro $C_2 \perp C_4$, $\sigma_h$, $2\sigma_v$, $2\sigma_d$, $i$, $S_4$:
$D_{4h}$. \ce{PCl5} (bipirâmide trigonal): $C_3$, três $C_2$, $\sigma_h$, $3\sigma_v$,
$S_3$: $D_{3h}$. 1,3,5-Triclorobenzeno: os mesmos elementos que \ce{BF3}: $D_{3h}$.
\end{solution}

\begin{solution}{exo:b3:point-groups:2}
Etano alternado: $C_3$, três $C_2$ perpendiculares, três $\sigma_d$, $i$, $S_6$:
$D_{3d}$. Eclipsado: $C_3$, três $C_2$, $\sigma_h$, três $\sigma_v$: $D_{3h}$.
Ciclo-hexano em cadeira: $D_{3d}$.
\end{solution}

\begin{solution}{exo:b3:point-groups:3}
\ce{SO3} ($D_{3h}$): não. \ce{SO2} ($C_{2v}$): sim. \ce{PCl3} ($C_{3v}$): sim.
\ce{XeF4} ($D_{4h}$): não. \ce{CH2Cl2} ($C_{2v}$): sim. cis-1,2-dicloroeteno
($C_{2v}$): sim; trans ($C_{2h}$): não.
\end{solution}

\begin{solution}{exo:b3:point-groups:4}
$C_2(z) = \mathrm{diag}(-1,-1,1)$, $\chi = -1$; $i = \mathrm{diag}(-1,-1,-1)$,
$\chi = -3$; $\sigma_h(xy) = \mathrm{diag}(1,1,-1)$, $\chi = 1$.
\end{solution}

\begin{solution}{exo:b3:point-groups:5}
Com as matrizes diagonais do exercício anterior: $C_2i = \sigma_h$, $C_2\sigma_h =
i$, $i\sigma_h = C_2$, cada operação é seu próprio inverso e todos os produtos
comutam. O grupo é comutativo, logo $X^{-1}AX = A$ para todo $X$: cada
operação forma sozinha uma classe (quatro classes, quatro \omterm{def:b3:point-groups:irrep}{representações
irredutíveis} de dimensão 1).
\end{solution}

\begin{solution}{exo:b3:point-groups:6}
Tome $\sigma_v(1)$ como o plano $xz$. O ponto $(1,0)$ vai por $C_3$ a $(-\frac12,
\frac{\sqrt3}2)$, depois por $\sigma_v(1)$ a $(-\frac12,-\frac{\sqrt3}2)$; na outra
ordem, vai a $(1,0)$ e depois a $(-\frac12,\frac{\sqrt3}2)$. Os dois produtos são
reflexões em dois planos diferentes ($\sigma_v(3)$ e $\sigma_v(2)$): o grupo não é
comutativo, e é por isso que suas reflexões formam uma classe de três.
\end{solution}

\begin{solution}{exo:b3:point-groups:7}
A bifenila torcida conserva o $C_2$ ao longo da ligação entre os anéis e dois $C_2$
perpendiculares a ele, que dividem ao meio os ângulos entre os planos dos anéis;
todos os planos e o \omterm{def:b3:point-groups:inversion}{centro de inversão} das formas plana ou perpendicular se
perdem. Com três $C_2$ perpendiculares e nenhum elemento impróprio, o grupo é
$D_2$: a molécula é quiral (suas duas formas torcidas são enantiômeros,
separáveis quando substituintes orto volumosos impedem a rotação).
\end{solution}

\begin{solution}{exo:b3:point-groups:8}
$d_{z^2}$ e $d_{x^2-y^2}$: $a_1$ ($x^2$, $y^2$, $z^2$ são $A_1$); $d_{xy}$: $a_2$; $d_{xz}$:
$b_1$; $d_{yz}$: $b_2$.
\end{solution}

\begin{solution}{exo:b3:point-groups:9}
$1 + 1 + 4 + 9 + 9 = 24 = h$. $\sum g\,\chi_{A_1}\chi_{T_2} = 1\cdot3 + 8\cdot0 + 3(-1)
+ 6(-1) + 6(1) = 0$.
\end{solution}

\begin{solution}{exo:b3:point-groups:10}
No plano perpendicular à reta de interseção, uma reflexão em uma reta
de ângulo $\alpha$ leva o ângulo polar $\phi$ a $2\alpha - \phi$. Duas reflexões, em
$\alpha$ e depois em $\alpha + \theta$: $\phi \to 2\alpha - \phi \to 2(\alpha + \theta) - (2\alpha - \phi) =
\phi + 2\theta$, uma rotação de $2\theta$ em torno da reta. Dois planos verticais a
$60^\circ$ geram, portanto, uma rotação de $120^\circ$: um $C_3$.
\end{solution}

\begin{solution}{exo:b3:point-groups:11}
O eixo C=C=C é um $C_2$ e um $S_4$ (um giro de $90^\circ$, que troca os dois
planos \ce{CH2}, seguido de uma reflexão no plano do carbono central); os dois
planos \ce{CH2} são $\sigma_d$; dois $C_2$ perpendiculares ao eixo os dividem ao meio.
Ordem 8: $D_{2d}$, aquiral. Em \ce{ClHC=C=CHCl}, os planos, o $S_4$ e o $C_2$
axial são destruídos pela substituição; só sobrevive um $C_2$ perpendicular ao
eixo: $C_2$, quiral --- um aleno com quiralidade axial.
\end{solution}

\begin{solution}{exo:b3:point-groups:12}
\ce{SF6}: $O_h$ ($h = 48$). \ce{SF5Cl}: o $C_4$ que passa por Cl e quatro $\sigma_v$:
$C_{4v}$ ($h = 8$). trans-\ce{SF4Cl2}: $D_{4h}$ ($h = 16$). cis-\ce{SF4Cl2}: $C_{2v}$ ($h
= 4$). Cada um conserva um subconjunto das operações de $O_h$, fechado para os
produtos, e 8, 16, 4 dividem 48. Polares: \ce{SF5Cl} e cis-\ce{SF4Cl2}.
\end{solution}

\begin{solution}{pb:b3:point-groups:1}
\textbf{1.} Ao longo das quatro diagonais do cubo que passam por C e por cada H;
cada uma dá $C_3$ e $C_3^2$: 8 rotações.
\textbf{2.} Ao longo de $x$, $y$, $z$, pelos centros das faces: $C_2(z)$ leva $(1,1,1)
\to (-1,-1,1)$, um hidrogênio. Uma rotação de $90^\circ$ em torno de $z$ seguida
de uma reflexão em $xy$ leva $(1,1,1) \to (-1,1,1) \to (-1,1,-1)$, um hidrogênio: um
$S_4$; com $S_4^3$, 6 operações.
\textbf{3.} Os planos $x = \pm y$, $y = \pm z$, $x = \pm z$; $x = y$ contém
$(1,1,1)$ e $(-1,-1,1)$: cada plano contém duas ligações C--H.
\textbf{4.} $1 + 8 + 3 + 6 + 6 = 24$.
\textbf{5.} A inversão leva $(1,1,1)$ a $(-1,-1,-1)$, que não é um hidrogênio:
não há $i$.
\textbf{6.} $E$; $8C_3$; $3C_2$; $6S_4$; $6\sigma_d$.
\textbf{7.} \ce{CH3Cl}: $C_{3v}$, $h = 6$.
\textbf{8.} \ce{CH2Cl2}: $C_{2v}$, $h = 4$.
\textbf{9.} \ce{CHCl3}: $C_{3v}$; \ce{CCl4}: $T_d$, $h = 24$.
\textbf{10.} \ce{CH2FCl}: $C_s$ ($h = 2$); \ce{CHFClBr}: $C_1$ ($h = 1$).
\textbf{11.} Cada um conserva as operações de $T_d$ que respeitam a
substituição; 6, 4, 6, 24, 2, 1 dividem todos 24.
\textbf{12.} Os três hidrogênios, permutados por $C_3$, $C_3^2$ e pelos três
$\sigma_v$.
\textbf{13.} Todas, exceto \ce{CCl4}: ao longo do eixo $C_3$ para \ce{CH3Cl} e \ce{CHCl3},
do eixo $C_2$ para \ce{CH2Cl2}, no \omterm{def:b3:point-groups:mirror}{plano especular} para \ce{CH2FCl}, em nenhuma
direção imposta para \ce{CHFClBr}.
\textbf{14.} Apenas \ce{CHFClBr} ($C_1$, sem $S_n$).
\textbf{15.} Dois enantiômeros.
\textbf{16.} Seus dois hidrogênios são equivalentes: o plano que passa por C, F
e Cl e divide ao meio H--C--H é um \omterm{def:b3:point-groups:mirror}{plano especular}.
\textbf{17.} Não: quatro substituintes diferentes deixam apenas $E$; sem $S_n$, a
molécula é quiral (e pode ser polar).
\textbf{18.} $(x, y, z)$: $T_2$.
\textbf{19.} Dois hidrogênios são as extremidades de uma aresta do tetraedro; as 24
operações levam qualquer aresta sobre qualquer outra (as 6 arestas formam um
só conjunto): um único \ce{CH2Cl2}.
\textbf{20.} \ce{CH2FCl}: um. \ce{CHFClBr}: dois, um par de enantiômeros.
\textbf{21.} Três: orto ($C_{2v}$), meta ($C_{2v}$), para ($D_{2h}$).
\textbf{22.} Três: 1,2,3- ($C_{2v}$), 1,2,4- ($C_s$), 1,3,5- ($D_{3h}$).
\textbf{23.} Dois (átomos de Cl adjacentes ou opostos). Só se conhece um
\ce{CH2Cl2}, o que excluiu um carbono plano e apoiou o carbono tetraédrico
proposto em 1874.
\textbf{24.} $1^2 + 1^2 + 2^2 + 3^2 + 3^2 = 24$: \textbf{a ordem de $T_d$, o
\omterm{def:b3:point-groups:point-group}{grupo pontual} do metano, é $h = 24$}.
\end{solution}
